\scanpage{158}
\uline{Catégorie galoisienne}

(i) Une catégorie $\unsure{\mathcal C|\mathscr C}$

(ii) Un foncteur $F:\mathcal C\to\text{Ensembles finis}$

Axiomes :

\margin{Poser plutôt en axiome l'\uline{existence} d'un $F$ ayant les \uuline{propriétés} : \deleted{\unsure{Soit|unicité}} $F$ permute aux sommes, aux produits, au \unsure{passage au quotient par un} groupe fini, et est tel qu'un \unsure{isomorphisme|injection} $F(X)\mathrel{\unsure{-|\to}} F(Y)$ \unsure{$X$ induit.|induit.|$X$ isom.} $\Rightarrow X\xrightarrow{\sim}Y$.}

(Gal 1) Dans $\mathcal C$, \deleted{\unsure{le|les}} sommes et \deleted{\unsure{le|les}} produits \deleted{de deux objets} \inserted{finis} existent
\inserted{[d'où l'existence de $\underline{0}$ et $\underline{1}$]}

\deleted{(Gal 3) $X\in\mathcal C$, $F(X)=\varnothing\Rightarrow X$ est une unité à gauche (soit $\underline{0}$)}

(Gal 4) $F$ permute à la somme et au produit

% margin note written sideways
\margin{Gal 3 et Gal 5 ensemble sont \unsure{impliqués} par : Si $f:X\to Y$ est tel que $F(X)\mathrel{\unsure{-|\to}}F(Y)$ soit \uline{injectif} \deleted{\unsure{($Y$ connexe)}} \inserted{($X$ \unsure{et|,} \deleted{\unsure{$Y$}} \unsure{connexe|connexes})} alors c'est un isomorphisme \inserted{de} $X$ sur une composante connexe \unsure{$Y_i$}.}

$X\in\mathcal C$ est dit \deleted{\uline{\unsure{irréductible|indivisible|indécomposable}}} \inserted{connexe} si $X=X'\amalg X''$ implique $X'=\underline{0}$
ou $X''=\underline{0}$

(Gal 2) $X\in\mathcal C$ \deleted{\unsure{$X\ne\underline{0}$}} implique $X\simeq\coprod_{i\in I}X_i$, \inserted{$X_i\ne\underline{0}$}, et tte décomposition
$X\simeq X'\amalg X''$ \unsure{impliques|implique} $X'\simeq\coprod_{i\in I'}X_i$, $X''=\coprod_{i\in I''}X_i$,
$I'\cap I''=\varnothing$, $I'\cup I''=I$. Les $X_i$, déterminés de
façon essentiellement unique [avec leur $X_i\to X$]
sont appelés les \uline{comp.} \deleted{\uline{\unsure{irréductibles|indivisibles|indécomposables}}} \inserted{connexes} de $X'$ [N.B. $X=\underline{0}\Leftrightarrow I=\varnothing$]

(\uuline{Gal 3}) Soit $f:X\to Y$ avec $Y$ \deleted{\unsure{irréductible|indivisible|indécomposable}} \inserted{connexe $\ne\underline{0}$}. Alors les
composantes \inserted{connexes} $X_i$ de $X$ telles que $f$ induise
un \uline{isom.} $f_i:X_i\xrightarrow{\sim}Y$ sont en correspondance
biunivoque avec les \uline{sections} $g_i:Y\to X$ \deleted{\unsure{$f$ de|de $f$}}
[hom. tels que $fg=\mathrm{id}_Y$], $g_i$ étant le composé
\deleted{\unsure{$Y\,X$}} $Y\xrightarrow{f_i^{-1}}X_i\xrightarrow{\text{inj. can}}X$ [en d'autres termes,
\deleted{toute} section se factorise] [N.B. il suffit
de le vérifier si $X=Z\times Y$, $f=p_2$]

(Gal 5) Sous les conditions de Gal 3, si \deleted{\unsure{$X$ est}} \inserted{$X,Y$ connexes et} $f:X\to Y$
induit $F(X)\xrightarrow{\sim}F(Y)$, c'est un isomorphisme
[implique $F(X)=\varnothing\to X=\underline{0}$] [Il suffit
de vérifier Gal 5 quand $X$ est une composante
connexe d'un produit $Z\times Y$, $f$ induit par $p_2$]

(\uuline{Gal 6}) Soit $X\in\mathcal C$, $X$ connexe $\ne\underline{0}$, et \uline{galoisien}, \deleted{\unsure{soit|i.e.}} i.e.
\deleted{\unsure{pour|Alors}} $G=\operatorname{Aut}(X)$ opère \uline{transitivement} dans $F(X)$
\deleted{\unsure{Tte composante connexe d'une puissance finie de} \illegible{1}}
Alors pour tout \unsure{ss-groupe} $H$ de $G$, $X/H$ dans $\mathcal C$ \uline{existe},
et $F(X/H)\simeq F(X)/H$
% archive sheet number: 157
